\section{Endpoint conjugation and central
twists}\label{endpoint-conjugation-and-central-twists}

Let \((G;a,b,(r_i)_{i\in E})\) be a C-group marking with E finite and \([a,r_i]=1\)
for all \(i\in E\). No string diagram, finite group, or underlying
geometry is assumed. Define

\[
a'=aba,\quad R_J=\langle r_j:j\in J\rangle,\quad
U_J=\langle a,b,R_J\rangle,\quad H_J=\langle a',b,R_J\rangle,
\] \[
P=\langle b,R_E\rangle,\quad H=H_E,\quad
\mathcal F=\{J\subseteq E:a\in H_J\}.
\tag{E1}
\]

Conjugation by a exchanges b and \(a'\) and fixes \(R_J\) pointwise.
Consequently

\[
H_J\trianglelefteq U_J=H_J\langle a\rangle,
\qquad [U_J:H_J]\le2.
\tag{E2}
\]

\textbf{Theorem C.Endpoint.} The new marking \((H;a',b,(r_i))\) has IP
if and only if

\[
P\cap P^a=R_E,
\qquad \mathcal F\text{ is closed under intersections}.
\tag{E3}
\]

The group H is normal of index at most two in G. It is G precisely when
\(\mathcal F\) is nonempty. Since E is finite, a nonempty
intersection-closed \(\mathcal F\) is precisely the family of supersets
of one least support \(S\subseteq E\).

\textbf{Proof.} The family \(\mathcal F\) is upward closed. Any new
marked subgroup is of one of the forms

\[
A_{\epsilon,J}=\langle b^\epsilon,R_J\rangle\ (\epsilon=0,1),
\qquad H_J,\qquad Q_J^a\text{ with }Q_J=\langle b,R_J\rangle,
\tag{E4}
\]

where \(b^{0}\) means that b is omitted. These correspond to omitting
\(a'\), using both endpoint labels, or using \(a'\) but omitting b.
Intersections of two A's are old marked intersections. Also

\[
A_{\epsilon,I}\cap H_J
\le A_{\epsilon,I}\cap U_J=A_{\epsilon,I\cap J};
\]

all generators on the right belong to \(H_J\), so equality holds.
Intersections of two \(Q^a\)'s follow by conjugating old IP. Conjugation
by a normalizes \(H_I\), so an H-with-\(Q^a\) intersection reduces to
the just proved H-with-A case. For an A-with-\(Q^a\) intersection, the
first condition in (E3) confines it to \(R_E\). Intersecting there and
using that a centralizes \(R_E\) gives \(R_{I\cap J}\).

It remains to intersect \(H_I\) and \(H_J\). Put \(K=I\cap J\). Old IP
gives

\[
H_K\le H_I\cap H_J\le U_I\cap U_J=U_K.
\tag{E5}
\]

If \(a\in H_K\) then \(H_K=U_K\). Otherwise \(H_K\) has index two in
\(U_K\), so the middle subgroup is either \(H_K\) or \(U_K\). The latter
case occurs exactly when a belongs to both \(H_I\) and \(H_J\), and
intersection closure excludes it. This proves all new marked
intersections. Necessity of the first condition follows by intersecting
the two endpoint-deleted new parabolics; necessity of the second follows
by intersecting \(H_I,H_J\). The generators are involutions and the
proved IP supplies independence. Equation (E2) with \(J=E\) gives the
group and index assertions. A finite nonempty upward intersection-closed
family contains its total intersection S, giving the least-support
description. \(\square\)

\textbf{Corollary C.Odd (odd endpoint and proper halving).} If
\(m=|ab|\) is odd, the only test in (E3) is the overlap
\(P\cap P^a=R_E\), and \(H=G\). If H is proper in G, the only test is
again the overlap, and \([G:H]=2\). In either case the endpoint product
has order m/gcd(m,2) for finite m; infinite m remains infinite.

Indeed \(a'b=(ab)^{2}\), and for odd m,

\[
a=(a'b)^{(m+1)/2}b\in H_\varnothing.
\]

Thus \(\mathcal F\) is respectively all subsets or empty. An even
endpoint order by itself does not imply that H is proper; the
least-support alternative in Theorem C.Endpoint cannot be discarded.

\textbf{Lemma C.Twist (central twist).} Let a generating family of a group G contain
involutions b,c with \(|bc|=m\) odd and \(m\ge 3\). Suppose
\(w\in G\) is a nonidentity involution centralizing b and c. Replacing
b by wb in that family preserves G and changes \(|bc|\) to \(2m\).

\textbf{Proof.} The commuting factors w and bc have coprime orders, so
\((wbc)^m=w\) and \(|wbc|=2m\). Hence the new generators recover w and
b. The new marker wb is a nonidentity involution: \(wb=1\) would make
\(b=w\) commute with c, contrary to \(m\ge 3\) odd. \(\square\)

A parabolic containing w,b,c is unchanged by this twist; one omitting b
is unchanged. Other intersections and typed Helly still require proof.
Lemma C.Twist is the auxiliary operation needed in the finite
application below, not a blanket geometric preservation theorem.
